\documentclass[11pt]{article}
\usepackage{amsmath,amssymb,amsthm}
\usepackage[margin=1in]{geometry}
\newtheorem{theorem}{Theorem}[section]
\newtheorem{lemma}[theorem]{Lemma}
\newtheorem{corollary}[theorem]{Corollary}
\theoremstyle{definition}
\newtheorem{definition}[theorem]{Definition}

\title{The oracle risk of linear shrinkage}
\author{}
\date{}

\begin{document}
\maketitle

\begin{abstract}
An observation $X$ with mean $\theta$ and variance $\sigma^2$ is shrunk towards zero by a factor $c$. We compute the
best factor and its risk under squared error, and compare that risk with the risks of the two simplest
estimators, $0$ and $X$.
\end{abstract}

\section{Setting}

Throughout, $\theta \ne 0$ and $\sigma > 0$ are real numbers.

\begin{definition}\label{def:risk}
For $c \in \mathbb{R}$, the \emph{risk} of the estimator $cX$ is
$R(c) = c^2 \sigma^2 + (1-c)^2 \theta^2$, its mean squared error when $X$ has mean $\theta$ and variance $\sigma^2$.
The \emph{oracle factor} is
$c^\ast = \theta^2 / (\theta^2 + \sigma^2)$.
The \emph{oracle risk} is
$\rho = \theta^2 \sigma^2 / (\theta^2 + \sigma^2)$.
It is the risk of the oracle factor, as Lemma~\ref{lem:square} shows.
\end{definition}

\section{Results}

\begin{lemma}\label{lem:square}
For every $c \in \mathbb{R}$, $R(c) - \rho = (\theta^2 + \sigma^2)(c - c^\ast)^2$.
\end{lemma}

\begin{proof}
Write $s = \theta^2 + \sigma^2$, so that $c^\ast = \theta^2 / s$ and $1 - c^\ast = \sigma^2 / s$. For
$c = c^\ast + d$,
\[
  R(c) = (c^\ast + d)^2 \sigma^2 + (1 - c^\ast - d)^2 \theta^2
       = \frac{\theta^4 \sigma^2 + \sigma^4 \theta^2}{s^2} + 2d\Bigl(\frac{\theta^2\sigma^2}{s} - \frac{\sigma^2\theta^2}{s}\Bigr) + s\, d^2
       = \rho + s\, d^2 ,
\]
which is the claim.
\end{proof}

\begin{lemma}\label{lem:min}
$\rho \le \min\{\theta^2, \sigma^2\}$.
\end{lemma}

\begin{proof}
Since $\theta^2 + \sigma^2 > \sigma^2 > 0$, we have
$\rho = \theta^2 \cdot \sigma^2 / (\theta^2 + \sigma^2) < \theta^2$, and in the same way $\rho < \sigma^2$.
\end{proof}

\begin{theorem}\label{thm:oracle}
For every $c \in \mathbb{R}$, $R(c) \ge \rho$, with equality if and only if $c = c^\ast$. Moreover,
$\rho \le \min\{R(0), R(1)\}$.
\end{theorem}

\begin{proof}
By Lemma~\ref{lem:square}, $R(c) - \rho = (\theta^2 + \sigma^2)(c - c^\ast)^2 \ge 0$, with equality if and only if
$c = c^\ast$, because $\theta^2 + \sigma^2 > 0$. Since $R(0) = \theta^2$ and $R(1) = \sigma^2$,
Lemma~\ref{lem:min} gives $\rho \le \min\{R(0), R(1)\}$.
\end{proof}

\begin{corollary}\label{cor:half}
$R(c^\ast) \le R(1/2)$.
\end{corollary}

\begin{proof}
By Lemma~\ref{lem:square}, $R(1/2) - R(c^\ast) = (\theta^2 + \sigma^2)(1/2 - c^\ast)^2 \ge 0$. By
Lemma~\ref{lem:min}, moreover, $R(c^\ast) = \rho \le \min\{\theta^2, \sigma^2\}$.
\end{proof}

\end{document}
